\documentclass[12pt,a4paper]{amsart}

\usepackage{amsmath,amssymb,amsthm}
\usepackage{mathtools}
\usepackage{hyperref}
\usepackage{geometry}
\geometry{margin=2.5cm}

% -------------------------------------------------------
% Theorem environments
% -------------------------------------------------------
\newtheorem{theorem}{Theorem}[section]
\newtheorem{lemma}[theorem]{Lemma}
\newtheorem{corollary}[theorem]{Corollary}
\newtheorem{proposition}[theorem]{Proposition}
\theoremstyle{definition}
\newtheorem{definition}[theorem]{Definition}
\newtheorem{remark}[theorem]{Remark}

% -------------------------------------------------------
% Custom commands
% -------------------------------------------------------
\newcommand{\N}{\mathbb{N}}
\newcommand{\Z}{\mathbb{Z}}
\newcommand{\R}{\mathbb{R}}
\newcommand{\C}{\mathbb{C}}
\DeclareMathOperator{\Re}{Re}
\DeclareMathOperator{\Li}{Li}

% -------------------------------------------------------
\begin{document}
% -------------------------------------------------------

\title{Riemann's Explicit Formula and the Prime Number Theorem}

\author{Michael Fuhrmann}
\address{Specialist Mathematics Project, Build 84}
\email{specialist-maths@localhost}
\date{\today}

\begin{abstract}
We develop Riemann's explicit formula connecting the prime-counting function
$\pi(x)$ to the zeros of $\zeta(s)$:
\[
  \pi(x) = \Li(x) - \sum_\rho \Li(x^\rho) - \log 2
  + \int_x^{\infty} \frac{dt}{t(t^2-1)\log t},
\]
where $\Li(x) = \int_2^x dt/\log t$ and the sum runs over all non-trivial
zeros $\rho$.  From this formula we derive the Prime Number Theorem
$\pi(x) \sim x/\log x$, and --- assuming the Riemann Hypothesis ---
the sharp error term
$|\pi(x) - \Li(x)| < \tfrac{1}{8\pi}\sqrt{x}\,\log x$ for $x \ge 2657$.
The von~Mangoldt function $\Lambda(n)$ and the Chebyshev functions
$\psi(x)$ and $\theta(x)$ are introduced and related to $\zeta$.
Finally, we indicate the connection to Dirichlet $L$-functions and GRH.
\end{abstract}

\maketitle
\tableofcontents

% ================================================================
\section{Introduction}
% ================================================================

The bridge between analytic properties of $\zeta(s)$ and arithmetic
properties of the primes is the \emph{explicit formula}.  Riemann sketched
it in his 1859 memoir \cite{Riemann1859}; it was made rigorous by von~Mangoldt
\cite{vonMangoldt1895} and Hadamard \cite{Hadamard1896}.

The fundamental idea: the location of the zeros of $\zeta(s)$ encodes
precise information about how primes are distributed.  If all zeros lay
on $\Re(s) = 1$ (which we know is false), we would get $\pi(x) \sim x/\log x$
with a weak error term.  The Riemann Hypothesis ($\Re(\rho) = 1/2$)
gives the sharpest possible error.

% ================================================================
\section{Arithmetic Functions and Chebyshev's Functions}
% ================================================================

\begin{definition}[Von~Mangoldt function]
\label{def:mangoldt}
The \emph{von~Mangoldt function} is defined by
\[
  \Lambda(n) \;=\;
  \begin{cases}
    \log p & \text{if } n = p^k \text{ for some prime } p \text{ and } k \ge 1, \\
    0      & \text{otherwise.}
  \end{cases}
\]
\end{definition}

\begin{definition}[Chebyshev functions]
\label{def:chebyshev}
The \emph{first Chebyshev function} is
$\theta(x) = \sum_{p \le x} \log p$,
and the \emph{second Chebyshev function} (or \emph{psi function}) is
\[
  \psi(x) = \sum_{n \le x} \Lambda(n) = \sum_{p^k \le x} \log p
           = \sum_{p \le x} \lfloor\log_p x\rfloor \cdot \log p.
\]
\end{definition}

\begin{proposition}[Relations between $\pi$, $\theta$, $\psi$]
\label{prop:psi_theta_pi}
\begin{align}
  \psi(x) &= \theta(x) + \theta(x^{1/2}) + \theta(x^{1/3}) + \cdots
  \label{eq:psi_theta} \\
  \pi(x) &= \frac{\theta(x)}{\log x} + \int_2^x \frac{\theta(t)}{t(\log t)^2}\,dt.
  \label{eq:pi_theta}
\end{align}
In particular, $\pi(x) \sim x/\log x$ is equivalent to $\psi(x) \sim x$
and to $\theta(x) \sim x$.
\end{proposition}

\begin{proof}
For~\eqref{eq:psi_theta}: each term $\theta(x^{1/k})$ contributes
$\sum_{p \le x^{1/k}} \log p$, accounting for prime powers $p^k \le x$.
For~\eqref{eq:pi_theta}: use Abel's summation formula
$\pi(x) = \theta(x)/\log x + \int_2^x \theta(t)/[t(\log t)^2]\,dt$,
which follows from $d[\pi(t)] = d[\theta(t)/\log t]$ and integration by parts.
\end{proof}

% ================================================================
\section{The Logarithmic Derivative and \texorpdfstring{$\psi$}{psi}}
% ================================================================

\begin{theorem}[Logarithmic derivative of $\zeta$]
\label{thm:log_deriv}
For $\Re(s) > 1$,
\[
  -\frac{\zeta'(s)}{\zeta(s)} \;=\; \sum_{n=1}^{\infty} \frac{\Lambda(n)}{n^s}.
\]
\end{theorem}

\begin{proof}
Taking the logarithm of the Euler product $\zeta(s) = \prod_p (1-p^{-s})^{-1}$:
\[
  \log\zeta(s) = -\sum_p \log(1 - p^{-s}) = \sum_p \sum_{k=1}^{\infty} \frac{p^{-ks}}{k}.
\]
Differentiating term by term (justified for $\Re(s) > 1$):
\[
  \frac{\zeta'(s)}{\zeta(s)} = -\sum_p \sum_{k=1}^{\infty} \frac{\log p}{p^{ks}}
  = -\sum_{n=1}^{\infty} \frac{\Lambda(n)}{n^s}. \qedhere
\]
\end{proof}

% ================================================================
\section{Perron's Formula and the Explicit Formula for \texorpdfstring{$\psi$}{psi}}
% ================================================================

\begin{theorem}[Perron's formula, simplified]
\label{thm:perron}
For $c > 1$ and $x > 0$ not a prime power,
\[
  \psi(x) = \frac{1}{2\pi i} \int_{c-i\infty}^{c+i\infty}
  \left(-\frac{\zeta'(s)}{\zeta(s)}\right) \frac{x^s}{s}\,ds.
\]
\end{theorem}

\begin{proof}[Proof sketch]
This follows from the Mellin transform and the inverse Laplace transform.
The key identity is: for $x > 0$, $c > 0$,
\[
  \frac{1}{2\pi i} \int_{c-i\infty}^{c+i\infty} \frac{x^s}{s}\,ds
  = \begin{cases} 1 & x > 1, \\ 0 & 0 < x < 1. \end{cases}
\]
Inserting the Dirichlet series for $-\zeta'/\zeta$ and integrating term by term
(justified by convergence estimates) gives the formula.
See \cite{Davenport2000}, Chapter~17.
\end{proof}

\begin{theorem}[Explicit formula for $\psi(x)$]
\label{thm:explicit_psi}
For $x > 1$ and $x$ not a prime power,
\begin{equation}
  \label{eq:explicit_psi}
  \psi(x) \;=\; x - \sum_\rho \frac{x^\rho}{\rho}
  - \frac{\zeta'(0)}{\zeta(0)} - \tfrac{1}{2}\log(1 - x^{-2}),
\end{equation}
where the sum runs over all non-trivial zeros $\rho$ of $\zeta$,
ordered as $|\Im(\rho)| \le T$ and then $T \to \infty$
(the sum converges conditionally in this symmetric ordering).

Here $\zeta'(0)/\zeta(0) = \log(2\pi) \approx 1.8379\ldots$\,.
\end{theorem}

\begin{proof}[Proof sketch]
Apply Perron's formula (Theorem~\ref{thm:perron}) and shift the contour
to the left, collecting residues at the poles and zeros of
$-\zeta'(s)/\zeta(s)$:
\begin{itemize}
  \item Residue at $s = 1$: $\zeta$ has a simple pole, so $-\zeta'/\zeta$
    has a simple pole with residue $-1$; contribution: $-x^1/1 \cdot (-1) = x$.
  \item Residues at non-trivial zeros $\rho$: each zero of $\zeta$ is a
    simple pole of $-\zeta'/\zeta$ with residue $1$; contribution: $-x^\rho/\rho$.
  \item Trivial zeros at $s = -2m$: contribution:
    $\sum_{m=1}^\infty x^{-2m}/(-2m) = \tfrac{1}{2}\log(1-x^{-2})$.
  \item Contribution of the pole of $\Gamma$ at $s = 0$: $-\zeta'(0)/\zeta(0)$.
\end{itemize}
A fully rigorous proof (handling convergence carefully) is in
\cite{Davenport2000}, Chapter~17.
\end{proof}

% ================================================================
\section{The Prime Number Theorem}
% ================================================================

\begin{theorem}[Prime Number Theorem]
\label{thm:pnt}
\[
  \pi(x) \;\sim\; \frac{x}{\log x} \quad (x \to \infty).
\]
Equivalently, $\psi(x) \sim x$.
\end{theorem}

\begin{proof}[Proof sketch via the explicit formula]
The key step is to show that $\zeta(s) \ne 0$ on the line $\Re(s) = 1$.

\textbf{Step 1} (No zeros on $\Re(s) = 1$):
Mertens (1898) observed: for $\sigma > 1$ and real $t$,
\[
  \log|\zeta^3(\sigma)\,\zeta^4(\sigma+it)\,\zeta(\sigma+2it)| \ge 0.
\]
This follows from $3 + 4\cos\theta + \cos 2\theta = 2(1+\cos\theta)^2 \ge 0$.
As $\sigma \to 1^+$, if $\zeta(1+it) = 0$, the fourth-order pole of
$\zeta^4$ would give $-\infty$ on the left, contradicting $\ge 0$.
Hence $\zeta(1+it) \ne 0$ for $t \ne 0$.

\textbf{Step 2} (PNT from zero-free region):
Once we know $\zeta(s) \ne 0$ on $\Re(s) = 1$, the explicit
formula~\eqref{eq:explicit_psi} shows the dominant term is $x$,
with all other terms $o(x)$.  Hence $\psi(x) \sim x$.
By Proposition~\ref{prop:psi_theta_pi}, $\pi(x) \sim x/\log x$.
\end{proof}

\begin{remark}
This approach is a streamlined version of the original proofs by
Hadamard \cite{Hadamard1896} and de~la~Vall\'ee~Poussin \cite{Poussin1896},
which established the PNT independently in 1896.
The elementary proof (without complex analysis) was found by
Erd\H{o}s and Selberg (1949).
\end{remark}

% ================================================================
\section{Error Terms: The Role of the Zeros}
% ================================================================

\begin{theorem}[Error term in terms of zeros]
\label{thm:error_term}
Under the assumption that the non-trivial zeros of $\zeta$ have real part
at most $\beta_0 < 1$, we have
\[
  \psi(x) = x + O\!\bigl(x^{\beta_0}(\log x)^2\bigr).
\]
In particular, the Riemann Hypothesis ($\beta_0 = \tfrac{1}{2}$) gives
\[
  \psi(x) = x + O(\sqrt{x}\,\log^2 x).
\]
\end{theorem}

\begin{proof}[Proof sketch]
In the explicit formula~\eqref{eq:explicit_psi}, the sum over zeros
$\sum_\rho x^\rho/\rho$ contributes at most
$\sum_\rho x^{\Re(\rho)}/|\rho| \le x^{\beta_0} \cdot \sum_\rho |\rho|^{-1}$.
The sum $\sum_\rho |\rho|^{-1+\epsilon}$ converges for $\epsilon > 0$
(since the number of zeros with $|\rho| \le T$ is $O(T\log T)$, so
$\sum_{|\rho|\le T} |\rho|^{-1} = O(\log^2 T)$).
\end{proof}

\begin{theorem}[Schoenfeld's explicit bound under RH]
\label{thm:schoenfeld}
Assuming the Riemann Hypothesis, for all $x \ge 2657$:
\[
  |\pi(x) - \Li(x)| \;<\; \frac{1}{8\pi}\,\sqrt{x}\,\log x.
\]
\end{theorem}

\begin{proof}[Proof sketch]
This follows from a careful analysis of the explicit formula and the
zero-free region under RH.  Schoenfeld \cite{Schoenfeld1976} established
this bound with explicit constants by carefully estimating the sum over
zeros $\sum_{|\gamma| \le T} x^{1/2+i\gamma}/(1/2+i\gamma)$ and the
tail $\sum_{|\gamma| > T}$.
\end{proof}

% ================================================================
\section{Riemann's Original Explicit Formula for \texorpdfstring{$\pi(x)$}{pi(x)}}
% ================================================================

\begin{theorem}[Riemann's explicit formula for $\pi(x)$]
\label{thm:riemann_explicit}
Define $\Li(x) = \int_2^x dt/\log t$.  Then for $x > 1$:
\begin{equation}
  \label{eq:riemann_explicit}
  \pi(x) \;=\; \Li(x) - \sum_\rho \Li(x^\rho) - \log 2
  + \int_x^{\infty} \frac{dt}{t(t^2-1)\log t}.
\end{equation}
Here the sum runs over all non-trivial zeros $\rho$ (ordered by
$|\Im(\rho)|$), and the integral (which converges for $x > 1$)
contributes $O(1/\log x)$.
\end{theorem}

\begin{proof}[Proof sketch]
This follows from the explicit formula for $\psi$ via the M\"obius
inversion formula.  Riemann used the ``prime-power counting function''
\[
  J(x) = \sum_{n=1}^{\infty} \frac{1}{n}\,\pi(x^{1/n})
  = \frac{1}{2\pi i}\int_{c-i\infty}^{c+i\infty} \frac{\log\zeta(s)}{s}\,x^s\,ds,
\]
from which $\pi(x)$ is recovered by M\"obius inversion:
$\pi(x) = \sum_{n=1}^\infty \mu(n)/n \cdot J(x^{1/n})$.
The poles and zeros of $\log\zeta(s)$ (via the Hadamard product)
then yield~\eqref{eq:riemann_explicit}.
A full treatment is in \cite{Titchmarsh1986}, Chapter~3.
\end{proof}

\begin{remark}
The integral in~\eqref{eq:riemann_explicit} equals $O\!\left((x\log x)^{-1}\right)$
as $x \to \infty$, and $\log 2 = 0.693\ldots$ is a constant.  The leading
term $\Li(x)$ dominates, and the correction $-\sum_\rho \Li(x^\rho)$ contains
the contribution of each non-trivial zero.
\end{remark}

% ================================================================
\section{Connection to \texorpdfstring{$L$}{L}-Functions and GRH}
% ================================================================

\begin{definition}[Dirichlet $L$-functions and primes in progressions]
\label{def:pnt_progressions}
For $\gcd(a, q) = 1$, let
\[
  \pi(x; q, a) = \#\{p \le x : p \equiv a \pmod q\}.
\]
By the orthogonality of Dirichlet characters,
\[
  \pi(x; q, a) = \frac{1}{\phi(q)}\sum_{\chi \bmod q} \bar\chi(a)\,
  \pi(x, \chi),
\]
where $\pi(x, \chi)$ involves zeros of $L(s, \chi)$.
\end{definition}

\begin{theorem}[Dirichlet's theorem with error term under GRH]
\label{thm:dirichlet_grh}
Assuming GRH for all $L(s, \chi)$ with $\chi \pmod q$,
\[
  \pi(x; q, a) = \frac{\Li(x)}{\phi(q)}
  + O\!\Bigl(\sqrt{x}\,\log(qx)\Bigr).
\]
In particular, primes are equidistributed in the $\phi(q)$ residue classes
coprime to $q$.
\end{theorem}

\begin{proof}[Proof sketch]
Apply the analogue of Theorem~\ref{thm:error_term} to each $L(s, \chi)$,
using $\Re(\rho_\chi) \le 1/2$ under GRH.  The character sum orthogonality
then gives the stated formula.
\end{proof}

% ================================================================
\section{Conclusion}
% ================================================================

Riemann's explicit formula is the central connection between the zeros
of $\zeta(s)$ and the distribution of prime numbers.
The Prime Number Theorem follows from the fact that $\zeta$ has no zeros
on the line $\Re(s) = 1$.
Under RH, the error term is $O(\sqrt{x}\log x)$,
essentially the sharpest possible estimate given the zero-free line hypothesis.
The formula~\eqref{eq:riemann_explicit} shows that each non-trivial zero
$\rho$ contributes an oscillation $\Li(x^\rho) \approx x^\rho/\log x$
to $\pi(x)$; the more zeros off the critical line, the larger these
oscillations would be.

% ================================================================
\begin{thebibliography}{10}

\bibitem{Riemann1859}
B.~Riemann,
\emph{\"Uber die Anzahl der Primzahlen unter einer gegebenen Gr\"o\ss e},
Monatsberichte der Berliner Akademie (1859), 671--680.

\bibitem{vonMangoldt1895}
H.~von~Mangoldt,
Zu Riemann's Abhandlung ``\"Uber die Anzahl der Primzahlen \ldots'',
\emph{J.\ Reine Angew.\ Math.}\ \textbf{114} (1895), 255--305.

\bibitem{Hadamard1896}
J.~Hadamard,
Sur la distribution des z\'eros de la fonction $\zeta(s)$ et ses
cons\'equences arithm\'etiques,
\emph{Bull.\ Soc.\ Math.\ France}\ \textbf{24} (1896), 199--220.

\bibitem{Poussin1896}
C.-J.\ de~la~Vall\'ee~Poussin,
Recherches analytiques sur la th\'eorie des nombres premiers,
\emph{Ann.\ Soc.\ Sci.\ Bruxelles}\ \textbf{20} (1896), 183--256.

\bibitem{Davenport2000}
H.~Davenport,
\emph{Multiplicative Number Theory}, 3rd~ed.,
Springer, New York, 2000.

\bibitem{Titchmarsh1986}
E.~C.~Titchmarsh,
\emph{The Theory of the Riemann Zeta-Function}, 2nd~ed.,
Oxford University Press, 1986.

\bibitem{Schoenfeld1976}
L.~Schoenfeld,
Sharper bounds for the Chebyshev functions $\theta(x)$ and $\psi(x)$, II,
\emph{Math.\ Comp.}\ \textbf{30} (1976), 337--360.

\bibitem{Apostol1976}
T.~M.~Apostol,
\emph{Introduction to Analytic Number Theory},
Springer, New York, 1976.

\end{thebibliography}

\end{document}
